% endtoend_fixedl.tex -- round 478, PRIME.RDAGGER.ENDTOEND_FIXEDL.01.
% PARTIAL(0.3-even-5x5-PD|tail-Schur).  NOT lambda_*(L)>=0.
% scramble-sensitive.  Companion: verify_endtoend_fixedl.py.
% Discovery: experiments/tfpt-discovery/endtoend_fixedl_probe.py.
% Fixed L only.  NO RH CLAIM.  NO anti-RH claim.
\documentclass[11pt]{article}

\usepackage[a4paper,margin=2.45cm]{geometry}
\usepackage{amsmath,amssymb,amsthm}
\usepackage{booktabs,array,microtype,xcolor}
\usepackage[T1]{fontenc}
\usepackage[colorlinks=true,linkcolor=blue!50!black,urlcolor=blue!50!black]{hyperref}

\newtheorem{lemma}{Lemma}
\newtheorem{theorem}{Theorem}
\theoremstyle{remark}
\newtheorem{remark}{Remark}
\newcommand{\QW}{Q_{\mathrm W}}
\setlength{\emergencystretch}{2em}

\title{\bfseries End-to-end fixed-$L$ Gram:\\[2mm]
the first block SATZ and the open tail\\[2mm]
\large PRIME.RDAGGER.ENDTOEND\_FIXEDL.01 --- round 478}
\author{Stefan Hamann}
\date{August 31, 2026}

\begin{document}
\maketitle

\begin{abstract}\noindent
At $L=0.3$ one has $P\equiv0$ and $t_0=22$.
The even Fourier Gram of $\QW$ on
$\mathrm{span}\{\cos(n\pi x/L):0\le n<N\}$
is interval-certified positive definite for $N=2,3,4,5$,
with every constant sealed.
The infinite tail Schur does not close:
the constant-row couplings decay like $1/n$ and the
Hilbert--Schmidt correction exceeds the Higham floor of the $3\times3$.
This is a SATZ on a finite-dimensional subspace, not
$\lambda_*(0.3)\ge0$ on $L^2[-L,L]$.
At $L=0.8$ the even Fourier $3\times3$ with exact rank-$3$ $P$
is indefinite; the best certified substructure remains the
r477 $2\times2$ of (constant, first Dirichlet).
The oscillatory multiplier
$\sigma_P(t)=\sum 2\Lambda(n)n^{-1/2}\cos(t\log n)$
is almost periodic: high-band gain versus $2S_{\mathrm{eff}}$
is at most $1.010$ at $L=0.8$ and does not unlock $L>1.6$.
Verdict \texttt{PARTIAL($0.3$-even-$5\times5$-PD$|$tail-Schur)}.
No RH claim.
\end{abstract}

\medskip\noindent
\textbf{Status.}\enspace
\texttt{PARTIAL($0.3$-even-$5\times5$-PD$|$tail-Schur)}.
No promotion of an RH claim.

\begin{center}
\fbox{\parbox{0.95\linewidth}{%
\textbf{Claim boundary.}  Nothing here is $\lambda_*(L)\ge0$
on all of $L^2[-L,L]$.  Finite even Fourier blocks are PD.
The tail Schur is open (explicit HS obstruction).
No silent class cut: the Paley--Wiener ideal
$\{\hat h=0\text{ on }|t|<t_0\}$ meets compact support only at $0$;
aliasing of interval Fourier modes is measured, not wished away.
No $L^*$ claim, no $R^\dagger$ claim, no RH claim,
no anti-RH claim.}}
\end{center}

\section{Convention}

After r461/r471, $\mathrm{supp}\,h\subset[-L,L]$,
$g=h*\widetilde h$, $\mathrm{supp}\,g\subset[-2L,2L]$,
$\QW=A-P+\Pi$ with
$P(g)=\sum_{n\le e^{2L}}2\Lambda(n)n^{-1/2}g(\log n)$
(scramble-sensitive) and $\Pi$ of rank at most two.
Plancherel as in r477:
$A(g)=(2\pi)^{-1}\int\sigma_A(t)|\hat h(t)|^2\,dt$,
$\sigma_A(t)=\mathrm{Re}\,\psi(\tfrac14+it/2)-\log\pi$.
Prime range $n\le e^{2L}$, not $e^L$.

\section{Even--odd split and the $t_0$-lemma}

\begin{lemma}[Even--odd at $P\equiv0$]\label{lem:split}
At $L=0.3$ there is no prime-power node with $\log n\le 0.6$,
so $P\equiv0$.  $A$ is a Fourier multiplier, hence
$A(\mathrm{even},\mathrm{odd})=0$.  The bilinear $\Pi$-density
of an odd $g$ vanishes.  Therefore
$\QW=Q_{\mathrm{even}}\oplus Q_{\mathrm{odd}}$
on $L^2[-L,L]$.
\end{lemma}

\begin{lemma}[$t_0$ and $c_\star$, r477]\label{lem:t0}
$t_0(0.3)=22$, $\sigma_A(22)_{\mathrm{lo}}=1.2517\ge1.2281
=4\sinh(0.3)+c_\star$ with $c_\star=0.01$, $S_{\mathrm{eff}}=0$.
On the ideal class $\{\hat h=0\text{ on }|t|<t_0\}$ one has
$\QW\ge c_\star\|h\|_2^2$.  That class intersects compact
support only at $h=0$.
\end{lemma}

\section{The block SATZ}

Write $u_0=1/\sqrt{2L}$ and
$u_n=\cos(n\pi x/L)/\sqrt{L}$ ($n\ge1$),
an ONB of even $L^2[-L,L]$.
The autocorrelation of $u_n$ ($n\ge1$) is
$g_n(t)=(L-t/2)\cos(\omega t)-\sin(\omega t)/(2\omega)$
with $\omega=n\pi/L$ (minus sine; the r476 first Dirichlet
uses a plus).  The $0$--$n$ correlation is the single
(not doubled) kernel
$g_{0n}(t)=(-1)^{n+1}\sin(\omega t)/\omega$.
$A+\Pi$ is the Bose series of r476/r477 with algebraic
$\psi$-tail and cubic pad, $K=200$.

\begin{theorem}[Even Fourier blocks at $L=0.3$]\label{thm:A}
Let $L=0.3$ and $G^{(N)}$ be the $N\times N$ Gram of $\QW$
in the basis $\{u_0,\dots,u_{N-1}\}$.
\begin{enumerate}
\item $N=2$: the interval characteristic polynomial gives
$\lambda_{\min}(G^{(2)})\in[1.731281\cdot10^{-2},\,1.825107\cdot10^{-2}]$.
Higham (raw) $\mu\ge8.301584\cdot10^{-3}$.
\item $N=3$: Higham (raw) $\mu\ge2.457362\cdot10^{-3}$;
midpoint $\lambda_{\min}\approx1.239183\cdot10^{-2}$.
\item $N=4,5$: r465 inverse-Cholesky congruence certifies
$G^{(N)}\succ0$; midpoint
$\lambda_{\min}\approx1.070798\cdot10^{-2}$ and
$1.027458\cdot10^{-2}$.
\end{enumerate}
The constant-mode Rayleigh is
$r(u_0)\in[1.58360545\cdot10^{-1},\,1.58360546\cdot10^{-1}]$.
The first odd mode $\sin(\pi x/L)$ has
$r\ge2.996973\cdot10^{-1}$.
\end{theorem}

\begin{remark}
The $3\times3$ Higham floor is conservative
(radius row-sum).  Midpoint eigenvalues decrease toward the
r476 first-Dirichlet Rayleigh
$r\in[1.163\cdot10^{-2},1.180\cdot10^{-2}]$, which lies
in the closed even span but not in any finite Fourier block.
\end{remark}

\section{Aliasing and the open tail}

Interval cosine $u_n$ is not bandlimited.
Low-frequency mass
$E_n=(2\pi\|u_n\|_2^2)^{-1}\int_{|t|<22}|\hat u_n|^2$:
$E_0\approx0.950$, $E_1\approx0.880$, $E_2\approx0.493$,
$E_3\approx0.0119$.  The low block $N\le3$ owns the band;
$n\ge3$ is already high.

\begin{lemma}[Tail Schur obstruction]\label{lem:tail}
The Hilbert--Schmidt mass of the couplings from the $3$-block
to modes $n=3,\dots,8$ satisfies
$\|G_{\ell h}\|_F\ge0.337$ (outward hulls).
The diagonal lower bound $r_3\ge1.678$ gives a Schur
correction $\ge6.783\cdot10^{-2}$, which exceeds the
Higham floor $2.457\cdot10^{-3}$ of $G^{(3)}$.
Gershgorin on row $0$ fails (harmonic $1/n$ couplings).
The infinite Gram is therefore not certified PD by this
Schur route.
\end{lemma}

$\Pi$ is rank one on the even sector
($\Pi(h)=2|\langle h,\cosh(x/2)\rangle|^2$) and is computed
exactly on every finite block; it is not the obstruction.
The obstruction is $A_{\ell h}$ (sinc overlap), not a crude
$4\sinh L$ majorant.

\section{$L=0.8$: cost and best substructure}

$t_0(0.8)=4300$, Paley--Wiener count
$N_{\mathrm{even}}\approx1095$, both parities $N\approx2190$.
Interval Cholesky of a $2190\times2190$ Bose Gram is
$\sim3.50\cdot10^9$ float flops plus $\sim 2\cdot10^6$
interval Bose entries: not a this-round computation.

\begin{theorem}[Best $L=0.8$ substructure]\label{thm:B}
The even Fourier $3\times3$ at $L=0.8$ with $P$ evaluated
exactly at the three nodes $n=2,3,4$
(linear functionals $g(\log n)$, not majorized)
is indefinite (midpoint $\lambda_{\min}\approx-4.72\cdot10^{-3}$).
The r477 Gram of
$\mathrm{span}\{1,\cos(\pi x/(2L))\}$ remains positive
definite: Schur$_{\mathrm{lo}}=7.413828\cdot10^{-2}$.
The constant-mode Rayleigh is $r(1)=8.027405\cdot10^{-2}$.
Not $\lambda_*(0.8)\ge0$.
\end{theorem}

\section{Oscillatory $P$}

As a Fourier multiplier,
$P=(2\pi)^{-1}\int\sigma_P(t)|\hat h(t)|^2\,dt$ with
\[
\sigma_P(t)=\sum_{n\le e^{2L}}2\Lambda(n)n^{-1/2}\cos(t\log n).
\]
The symbol is a finite cosine sum, hence almost periodic.
Its essential supremum on $[t_0,\infty)$ equals (up to an
arbitrarily small gap) its essential supremum on $\mathbb R$:
there is no $1/t_0$ decay of $\|\sigma_P\|_\infty$.

\begin{lemma}[Measured gain, scramble-sensitive]\label{lem:C}
Triangle: $|\sigma_P|\le 2S_{\mathrm{eff}}$.
Sampling $[t_0,t_0+200]$ gives an \emph{upper} bound on the
gain $2S_{\mathrm{eff}}/\|\sigma_P\|_\infty$:
$L=0.8$ (3 nodes) $\mathrm{gain}\le1.0034$;
$L=1.2$ (8 nodes) $\mathrm{gain}\le1.133$;
$L=1.8$ (18 nodes) $\mathrm{gain}\le1.836$.
A van der Corput saving for a single highly oscillatory $g$
is a different pairing and is not an operator-norm bound.
\end{lemma}

\section{Updated forecast}

The r477 crude law
$t_0\sim2\pi\exp(2S_{\mathrm{eff}}+4\sinh L)$
is unchanged in order.
Replacing $2S_{\mathrm{eff}}$ by the measured
$\|\sigma_P\|_\infty$ multiplies the exponent by
$1/\mathrm{gain}\approx1$ at $L=0.8$ and $\approx0.54$ at $L=1.8$;
$t_0(1.8)$ drops from $\sim10^{13}$ to still $\sim10^{10}$.
$L_{16}$ remains $\sim10^{39}$.  $L>1.6$ is not unlocked.

\section{Sealed balance}

Probe \texttt{endtoend\_fixedl\_probe.py},
SPEC\_SHA \texttt{138cbf9410486418}.
Smoke $17/17$ in $0.31\,\mathrm{s}$; full $27/27$ in $3.9\,\mathrm{s}$.
Two smoke runs pin-identical (elapsed excluded).

\begin{center}
\small
\begin{tabular}{@{}llr@{}}
\toprule
object & type & sealed constant\\
\midrule
$t_0(0.3)$ & r477 lemma & $22$\\
$G^{(2)}$ $\lambda_{\min}$ & interval charpoly & $[1.731\cdot10^{-2},1.826\cdot10^{-2}]$\\
$G^{(2)}$ Higham $\mu$ & raw & $\ge8.302\cdot10^{-3}$\\
$G^{(3)}$ Higham $\mu$ & raw & $\ge2.457\cdot10^{-3}$\\
$G^{(4)},G^{(5)}$ & congruence PD & midpoint $1.071,1.027\cdot10^{-2}$\\
tail HS correction & obstruction & $\ge6.783\cdot10^{-2}$\\
$L=0.8$ Fourier $3\times3$ & indefinite & hint $-4.72\cdot10^{-3}$\\
$L=0.8$ Dir.\ $2\times2$ & r477 Schur$_{\mathrm{lo}}$ & $7.414\cdot10^{-2}$\\
$K_P$ gain $L=0.8$ & $\le$ triangle & $\le1.010$\\
$N(0.8)$ both parities & cost & $2190$, $\sim3.50\cdot10^9$ flops\\
\bottomrule
\end{tabular}
\end{center}

\section{Kill table}

\begin{center}
\small
\begin{tabular}{@{}lll@{}}
\toprule
gate & status & note\\
\midrule
A-constants sealed & PASS & every number above is a pin\\
aliasing honest & PASS & $E_n$ measured; no class cut\\
tail Schur & KILL & HS $>$ Higham3; no $\lambda_*$\\
B full $2190$ Chol & KILL (cost) & forecast recorded\\
B Fourier $3\times3$ & KILL (indefinite) & exact $P$, not majorized\\
C $1/t_0$ gain & KILL & almost-periodic multiplier\\
scramble & PASS & literal $\log n$ in $P$ and $\sigma_P$\\
determinism & PASS & two-run pins identical\\
anti-list / no RH & PASS & fixed $L$, $e^{2L}$, no zeros\\
\bottomrule
\end{tabular}
\end{center}

\section{What this does not do}

It does not prove $\lambda_*(0.3)\ge0$ or $\lambda_*(0.8)\ge0$.
It does not close the Paley--Wiener remainder as a complete-space
argument.  It does not unlock $L>1.6$.  No cofinal quantifier.
No RH claim.

\end{document}
